\uuid{LtKv}
\titre{Triangularisation et décomposition nilpotente}
\niveau{L2}
\module{Algèbre}
\chapitre{Réduction d'endomorphisme, polynôme annulateur}
\sousChapitre{Trigonalisation}
\theme{Changement de base, matrice de passage, décomposition $D+N$ nilpotente}
\auteur{}
\datecreate{2026-06-04}
\organisation{AMSCC}
\difficulte{3}

\contenu{
\texte{
  Soit $A$ la matrice
  \[
    A = \begin{pmatrix}1&1&0\\-1&2&1\\-1&0&3\end{pmatrix}.
  \]
  On pose $P$ la matrice
  \[
    P = \begin{pmatrix}1&0&0\\1&1&0\\1&1&1\end{pmatrix}.
  \]
}
\begin{enumerate}

\item
  \question{Vérifier que $P^{-1}$ est la matrice
    \[
      \begin{pmatrix}1&0&0\\-1&1&0\\0&-1&1\end{pmatrix}.
    \]}
  \reponse{
    Il suffit de vérifier que $PP^{-1} = I_3$ :
    \[
      \begin{pmatrix}1&0&0\\1&1&0\\1&1&1\end{pmatrix}
      \begin{pmatrix}1&0&0\\-1&1&0\\0&-1&1\end{pmatrix}
      =
      \begin{pmatrix}
        1 & 0 & 0 \\
        1-1 & 1 & 0 \\
        1-1 & -1+1 & 1
      \end{pmatrix}
      = I_3. \quad\checkmark
    \]
  }

\item
  \question{Déterminer la matrice $T = P^{-1}AP$, puis écrire $T$ sous la forme
    $T = D + N$ où $D$ est une matrice diagonale et $N$ une matrice vérifiant $N^3 = 0$.}
  \reponse{
    \textbf{Calcul de $T = P^{-1}AP$.}

    On calcule d'abord $P^{-1}A$ :
    \[
      P^{-1}A
      = \begin{pmatrix}1&0&0\\-1&1&0\\0&-1&1\end{pmatrix}
        \begin{pmatrix}1&1&0\\-1&2&1\\-1&0&3\end{pmatrix}
      = \begin{pmatrix}1&1&0\\-2&1&1\\0&-2&2\end{pmatrix}.
    \]
    Puis $(P^{-1}A)P$ :
    \[
      T = \begin{pmatrix}1&1&0\\-2&1&1\\0&-2&2\end{pmatrix}
          \begin{pmatrix}1&0&0\\1&1&0\\1&1&1\end{pmatrix}
        = \begin{pmatrix}2&1&0\\0&2&1\\0&0&2\end{pmatrix}.
    \]

    \textbf{Décomposition $T = D + N$.}

    On pose :
    \[
      D = \begin{pmatrix}2&0&0\\0&2&0\\0&0&2\end{pmatrix} = 2I_3,
      \qquad
      N = \begin{pmatrix}0&1&0\\0&0&1\\0&0&0\end{pmatrix}.
    \]
    On vérifie bien que $T = D + N$. De plus :
    \[
      N^2 = \begin{pmatrix}0&0&1\\0&0&0\\0&0&0\end{pmatrix},
      \qquad
      N^3 = 0. \quad\checkmark
    \]
  }

\end{enumerate}
}
